\section{서론}

한국과학교육학회지(JKASE)는 과학교육 전반의 연구를 다루므로, 서론에서는 연구가 어떤 과학교육 문제와 연결되는지 명확히 제시하세요. 과학 교수학습, 교육과정, 평가, 교사교육, 과학영재교육, 과학문화, 교육정책 등 원고의 초점을 독자가 초반에 파악할 수 있도록 작성합니다.

선행연구 검토에서는 단순한 나열보다 연구 공백을 드러내는 흐름이 중요합니다. 국내외 과학교육 연구와 본 연구의 차별성을 연결해 주세요.

\subsection{학술지 및 투고 성격 확인}

한국과학교육학회지는 과학교육학자, 과학교육자, 과학교육 관련 전문가의 연구 결과를 공유하기 위한 학술지입니다. 투고 전 최신 투고규정에서 원고 유형, 초록 요구사항, 참고문헌 양식, 온라인 투고 절차를 확인하세요. 연구논문 외 논단, 논평, 해설, 소개 유형은 요구되는 글의 구조가 다를 수 있습니다.

\subsection{연구 목적}

연구 목적 또는 연구 문제를 입력하세요. 가능한 경우 연구 질문을 번호로 제시하면 심사자가 연구 범위와 분석 초점을 빠르게 확인할 수 있습니다.
